\documentclass[11pt,a4paper]{article}

% Fix fi/fl ligature extraction from the PDF (correct ToUnicode CMaps)
\input{glyphtounicode}
\pdfgentounicode=1
\usepackage{cmap}
\usepackage[utf8]{inputenc}
\usepackage[T1]{fontenc}
\usepackage{lmodern}
\usepackage{amsmath,amssymb,amsthm}
\usepackage{mathtools}
\usepackage{booktabs}
\usepackage{hyperref}
\usepackage{geometry}
\usepackage{enumitem}
\usepackage{listings}
\usepackage{longtable}
\usepackage{array}
\usepackage{xcolor}

\geometry{margin=1in}
\setcounter{tocdepth}{2}

\newtheorem{definition}{Definition}[section]
\newtheorem{principle}[definition]{Principle}
\newtheorem{theorem}[definition]{Theorem}
\newtheorem{proposition}[definition]{Proposition}
\newtheorem{invariant}[definition]{Invariant}
\newtheorem{remark}[definition]{Remark}

\lstset{
  basicstyle=\ttfamily\small,
  columns=fullflexible,
  keepspaces=true,
  xleftmargin=2em,
  aboveskip=0.5em,
  belowskip=0.5em,
  breaklines=true,
  showstringspaces=false
}

\title{The Ledger: The Managed-Account Workflow\\[6pt]
\large A derivation of \S 6 (Managed Accounts, Virtual Portfolios, and TRS)\\
from the primitives, with Addendum~A1 (StatesHome)\\
\normalsize Convergence of a nine-lens adversarial review}
\author{R.\ Delloye}
\date{June 2026 --- Workflow note to v10.3}

\begin{document}
\maketitle

\begin{abstract}
A managed account is not a primitive. It is the composition of four that already exist: a
wallet partition that holds positions and produces PnL; a mandate unit $u_{\mathrm{MA}}$
issued by the manager and held by the client, conserved by the issuance law; a
deterministic reset contract that observes performance, crystallises one net cash flow, and
resets its baseline; and the three-map state model of Addendum~A1, which homes the mandate's
terms, the shared observables, and the per-client economic state. This note derives the exact
managed-account workflow from those primitives and carries one numerical example
through subscription, a trade, a period's NAV, a performance-fee crystallisation against a
high-water mark, and redemption, verifying $\sum_w w(u)=0$ at every step. Nine specialist lenses
attacked the workflow with $140$ challenge questions; $105$ were answered and cleared, and $35$
escalated to three framework-level weaknesses that \S 6 cannot resolve without changing the
framework. The escalations are stated plainly in \S\ref{sec:escalations}; they are the most
valuable output of the review, not its defects.
\end{abstract}

\tableofcontents
\clearpage

%==============================================================================
\section{Primitives}
\label{sec:primitives}
%==============================================================================

\noindent\textbf{Wallet, unit, move, conservation.} A wallet $w\in\mathcal{W}$ has state
$w_t:\mathcal{U}\to\mathbb{R}$; $w_t(u)$ is the signed quantity of unit $u$, negative for a
short or obligation. A move $m=(w_s,w_d,u,q,t,\text{source},\text{meta})$ with $q>0$ acts by
$w_s(u)\mathbin{-}\!=q;\ w_d(u)\mathbin{+}\!=q$. A transaction is a finite set of moves applied
atomically. Conservation is algebraic, not a check:
\[
Q(u)=\sum_{w\in\mathcal{W}}w_t(u)=0\qquad\forall u,\ \forall t. \tag{P1}
\]
External counterparties are virtual wallets inside the closed system; the ledger has no
``outside'' within scope. A contract unit issued by party $A$ to party $B$ obeys the
\emph{issuance law} $w_A(u)=-1,\ w_B(u)=+1$, so $\sum_w w(u)=0$ by the same conservation.

\noindent\textbf{Valuation.} $V_t=\sum_{u}w_t(u)\,P_t(u)$ with $P_t$ an external price input
($P_t(\text{USD})=1$). Conservation guarantees \emph{quantity} integrity only; value moves with
$P_t$. PnL is path-independent: $\mathrm{PnL}=V_{t_1}-V_{t_0}$, and decomposes as
$\mathrm{PnL}=\mathrm{PnL}_{\text{price}}+\mathrm{PnL}_{\text{flow}}$. Arithmetic is
fixed-precision decimal with bankers' rounding only at instruction generation, so outputs are
bit-identical and replayable.

\noindent\textbf{The three-map state model (Addendum~A1).} State attaches to three maps:
immutable versioned \textit{ProductTerms}$[u]$; the shared \textit{UnitStatus}$[u]$, a
materialised projection of the log --- a cached view the log can always rebuild; and
per-position \textit{PositionState}$[w,u]$, with an \texttt{Option} accessor (\texttt{None}
``never held'' $\neq$ \texttt{Some(zero)} ``held, now flat'') and a monotone carrier (closed-out
rows are retained at zero). UnitStatus holds one value per unit, shared --- read identically by
every holder. Its value changes over time, but UnitStatus is not a separate source of truth: the
immutable event log is. Every change is caused by a logged event; the stored value is overwritten
in place only as a read cache. Replaying the events up to any point rebuilds the exact value that
held then, so nothing is lost by overwriting and there is no other way to change the value. A value
exists from registration onward. There is no wallet-keyed economic-state sector: every per-wallet
economic fact is $(w,u_{\mathrm{MA}})$-keyed. Condition~C2 requires each event handler to emit a
\texttt{StateDelta} that is structurally zero-sum per event class; C3 makes it atomic across all
three maps; C4 scopes reads by capability; C11 names the unique handler permitted to mutate each
\textit{PositionState} field.

%==============================================================================
\section{Account Establishment}
\label{sec:establishment}
%==============================================================================

\noindent\textbf{The mandate is a unit.} A managed account is opened by issuing the mandate unit
$u_{\mathrm{MA}}$ from the manager $w_M$ to the client $w_C$:
\begin{lstlisting}
Move(from: w_M, to: w_C, unit: u_MA, quantity: 1,
     source: "mandate_issue", metadata: "MANDATE_ISSUANCE")
\end{lstlisting}
giving $w_M(u_{\mathrm{MA}})=-1,\ w_C(u_{\mathrm{MA}})=+1,\ \sum_w w(u_{\mathrm{MA}})=0$. This is
a real two-named-wallet issuance with a real issuer and holder, not the rejected self-reflexive
sentinel; conservation holds by the issuance law (P1).

\begin{invariant}[Mandate is a singleton bilateral contract]
\label{inv:singleton}
$\operatorname{support}\big(w\mapsto w(u_{\mathrm{MA}})\big)=\{(w_M,-1),(w_C,+1)\}$ while the
mandate is live.
\end{invariant}

\noindent Conservation alone does \emph{not} establish Invariant~\ref{inv:singleton}: $P1$ admits
$(-5,+5)$ or a three-way split. Identity is pinned by four existing primitives together:
Invariant~\ref{inv:singleton} for the active phase; the \texttt{Option} accessor distinguishing
\texttt{None} (pre-issuance) from \texttt{Some(zero)} (terminated, retained); a
\textit{UnitStatus} lifecycle guard admitting issuance only from an unissued status (so a
terminated $u_{\mathrm{MA}}$ cannot be resurrected); and C8 routing a re-mandate to a fresh
$u'_{\mathrm{MA}}$ with \texttt{SupersededBy}.

\noindent\textbf{What is recorded where.}

\begin{longtable}{p{8.4cm}p{6.4cm}}
\toprule
\textbf{Datum} & \textbf{Home} \\
\midrule
Mandate text, fee schedule, benchmark identity, position limits, HWM/hurdle methodology,
crystallisation frequency & \textit{ProductTerms}$[u_{\mathrm{MA}}]$ (immutable, versioned) \\
Lifecycle stage, \texttt{superseded\_by} & \textit{UnitStatus}$[u_{\mathrm{MA}}]$ (shared) \\
HWM value, entry NAV, accrued fees, mandate-breach flags, subscription/redemption cursor,
benchmark NAV at this wallet's inception & \textit{PositionState}$[w_C,u_{\mathrm{MA}}]$ \\
Current benchmark level (an externally-sourced value, captured as a logged observation event) & \textit{UnitStatus}$[u_{\mathrm{bench}}]$ \\
\bottomrule
\end{longtable}

\noindent The placement of the per-client \emph{economic} scalars is the subject of
escalation~E1 (\S\ref{sec:escalations}): A1 stores them, yet each is a fold of the event log.

\begin{principle}[Mandate is non-valued]
\label{prin:nonvalued}
$u_{\mathrm{MA}}$ is excluded from the domain of valuation: $P_t(u_{\mathrm{MA}})$ is undefined,
not zero. A price on $u_{\mathrm{MA}}$ would add $+1\cdot P_t(u_{\mathrm{MA}})$ to the client's
per-wallet NAV and double-count the underlying exposure; the term cancels only in the global sum.
\end{principle}

%==============================================================================
\section{Subscription}
\label{sec:subscription}
%==============================================================================

Capital enters the managed wallet as a cash move from the client's external funding wallet
$w_{\text{ext}}$ (a virtual wallet at the ledger boundary):
\begin{lstlisting}
Move(from: w_ext, to: w_C, unit: USD, quantity: K,
     source: "subscribe", metadata: "SUBSCRIPTION")
\end{lstlisting}
The move is \textbf{tagged} \texttt{SUBSCRIPTION} and indexed by the subscription/redemption
cursor in \textit{PositionState}$[w_C,u_{\mathrm{MA}}]$. The tag is load-bearing: it lets the net
external flow over any interval be reconstructed as a projection over the stream
(\S\ref{sec:fees}). At subscription the entry NAV and the initial high-water mark are set to the
subscribed capital base; the high-water mark before subscription is \texttt{None}
(unrepresentable), never $0$ --- a $0$ mark would charge a performance fee on the entire
contributed capital at the first crystallisation. The benchmark NAV at inception is recorded from
\textit{UnitStatus}$[u_{\mathrm{bench}}]$ for performance \emph{reporting}.

%==============================================================================
\section{Trading under Mandate}
\label{sec:trading}
%==============================================================================

A trade is an ordinary two-legged transaction against a virtual market wallet $w_{\text{mkt}}$,
admitted only if it satisfies the mandate guards.

\begin{definition}[Mandate guard]
A quantitative mandate constraint is a precondition $g$ on the move-generation function. The
contract evaluates $g$ against the current state and the recorded price snapshot; if $g$ fails,
\emph{no} moves are emitted and the state is unchanged.
\end{definition}

\noindent Guard determinism is \textbf{price-relative}. The claim ``same wallet state $+$ same
trade $\Rightarrow$ same decision'' is false for value-based limits (leverage,
concentration-by-value, currency caps): holding positions fixed and doubling a price flips the
same trade from admit to reject. Determinism is recovered only when the guard is a pure function
of $(\textit{ProductTerms},\,\text{state},\,\text{trade},\,P_t)$ with $P_t$ an \emph{explicit,
snapshot-pinned} input. Two totality obligations follow: a missing price
($P_t(u)=\texttt{None}$) must make the guard \emph{fail closed} (a typed reject), never a silent
admit; and a \emph{passive} breach --- a limit crossed by price appreciation with no admitted
trade --- lies outside the domain of a move-time guard and is detected only by a periodic
valuation sweep, not by the guard. The pinned-snapshot requirement is escalation~E2.

%==============================================================================
\section{Valuation and NAV}
\label{sec:nav}
%==============================================================================

The account NAV is the wallet valuation with the mandate excluded
(Principle~\ref{prin:nonvalued}):
\[
\mathrm{NAV}_t \;=\; \sum_{u\neq u_{\mathrm{MA}}} w_{C,t}(u)\,P_t(u).
\]
By state-sufficiency, $\mathrm{NAV}_t$ depends only on current balances, current unit state, and
current prices --- not on history. Performance against the benchmark is computable from the
inception benchmark NAV and the current level in \textit{UnitStatus}$[u_{\mathrm{bench}}]$; it
feeds \emph{reporting}, and feeds the \emph{fee} only when the mandate's methodology names a
benchmark hurdle.

%==============================================================================
\section{Fee Logic}
\label{sec:fees}
%==============================================================================

\noindent\textbf{Performance is net of capital flows.} Over a reset interval
$[t_{k-1},t_k]$ the performance is the gross NAV change less the net external flow:
\[
\mathrm{Perf}_k \;=\; \big(\mathrm{NAV}_{t_k}-\mathrm{NAV}_{t_{k-1}}\big)\;-\;
\mathrm{NetFlow}_{[t_{k-1},t_k]},\qquad
\mathrm{NetFlow}=\!\!\sum_{\substack{m\ \text{tagged sub/redemption}\\ t_{k-1}<m.t\le t_k}}\!\!\pm\,m.q.
\]
The gross change $\mathrm{NAV}_{t_k}-\mathrm{NAV}_{t_{k-1}}$ conflates performance with capital
flows: by the PnL decomposition a mid-period subscription is a \emph{flow}, and crystallising it
as performance would charge a fee on the client's own contributed capital. $\mathrm{NetFlow}$ is a
projection over the already-tagged cursor moves (\S\ref{sec:subscription}); the correction adds no
new coordinate.

\noindent\textbf{Fee is not PnL settlement.} A single whole-$\mathrm{Perf}\to$beneficiary move is
the \emph{internal desk} case, where a book is swept to Treasury. A client managed account pays
the manager only a fee fraction; the gain \emph{stays in the client book}. The fee handler is
total over $\operatorname{sign}(\mathrm{NAV}-\mathrm{HWM})$:
\begin{align*}
\text{management fee } f^{m}_k &= r_{\text{m}}\cdot\Delta t_k\cdot \mathrm{NAV}^{\text{base}}_k
&&\text{(accrues on AUM regardless of sign),}\\
\text{performance fee } f^{p}_k &= r_{\text{p}}\cdot\max\!\big(0,\ \mathrm{NAV}^{\text{net}}_k-\mathrm{HWM}_{k-1}\big)
&&\text{(floored at zero; no clawback),}\\
\mathrm{HWM}_k &= \max\!\big(\mathrm{HWM}_{k-1},\ \mathrm{NAV}^{\text{net}}_k-f^{p}_k\big)
&&\text{(monotone ratchet; C11 writer \texttt{fee\_crystallise}),}
\end{align*}
where $\mathrm{NAV}^{\text{net}}_k=\mathrm{NAV}_{t_k}-f^{m}_k$. On a loss period the performance
fee is $0$ and the high-water mark is unchanged, while the management fee still accrues.

\noindent\textbf{Crystallisation is double-entry.} A fee is realised as a real cash move
$w_C\to w_M$, so it is conserved and carries its own contra-entry; a non-conserved stored scalar
would be single-entry and sit outside P1. The handler emits, atomically (C3):
\begin{lstlisting}
Move(from: w_C, to: w_M, unit: USD, quantity: f_m_k, metadata: "MGMT_FEE")
Move(from: w_C, to: w_M, unit: USD, quantity: f_p_k, metadata: "PERF_FEE")
\end{lstlisting}
\begin{theorem}[Fee zero-sum per crystallisation]
$\sum_w \Delta_{\text{fee}}(\text{USD})=0$ at each crystallisation: the client pays exactly what
the manager receives, creating and destroying no value.
\end{theorem}
\noindent\textit{Proof.} Each leg is $w_C\mathbin{-}\!=q;\ w_M\mathbin{+}\!=q$; summing over
wallets gives $-q+q=0$ per leg, hence $0$ over the transaction. \qed

\noindent\textbf{Quantity, baseline, residual.} The crystallise leg, like every settlement leg in
this note, emits $q=|x|$ with direction fixed by $\operatorname{sign}(x)$ and \emph{no} move at
$x=0$; a fixed-direction listing with a signed quantity is non-representable on a loss and illegal
at $q=0$. The reset baseline is the \emph{post-settlement} capital base
$B_k=\mathrm{NAV}_{t_k}-(f^{m}_k+f^{p}_k)+\mathrm{NetFlow}$; reading the pre-settlement NAV would
claw the distributed performance back next period. The rounding residual
$\mathrm{Perf}_k-\operatorname{round}(\mathrm{Perf}_k)$ is \emph{retained} in $w_C$ as a conserved
un-crystallised remainder, never dropped; it is the reconciling item between economic performance
and moved cash. Whether $B_k$, $\mathrm{HWM}_k$, and the
accrued fee are stored or re-derived is escalation~E1.

%==============================================================================
\section{Segregation}
\label{sec:segregation}
%==============================================================================

\begin{theorem}[Segregation]
\label{thm:seg}
No conservative move can perturb a wallet it does not name. Client partitions are mutually
isolated \emph{iff} conservation is combined with capability scoping.
\end{theorem}

\noindent\textbf{Conservation alone is not segregation.} A move
$m=(w_{A},w_{B},\text{USD},q)$ between two \emph{different} clients' partitions is perfectly
conservative ($\sum_w\Delta=0$), so $P1$ does not forbid it. The \S 6 phrasing ``conservation
enforces segregation by algebra'' conflates two independent facts: conservation (CONS,
$\sum_w\Delta=0$) and locality (LOC, $\operatorname{supp}(\Delta)\subseteq\{w_s,w_d\}$). Locality
gives ``no effect beyond the named wallets''; isolation of \emph{distinct} clients additionally
requires an authorisation predicate --- C4 capability scoping, forbidding a handler from naming or
reading across a $(w,u_{\mathrm{MA}})$ boundary it has no capability for. Segregation is therefore
a theorem under CONS $\wedge$ LOC $\wedge$ C4, which supplies the logical-segregation component of
CASS~6 / MiFID~II~16(8). Legal segregation (distinct custodial accounts, trust/nominee) is outside
scope: the ledger evidences it, it does not establish it. That C4 is asserted in prose rather than
typed on the accessor signature is escalation~E4 (\S\ref{sec:escalations}).

%==============================================================================
\section{CSA Margin and Collateral}
\label{sec:csa}
%==============================================================================

The CSA margin contract is a wallet-level smart contract attached to a per-counterparty collateral
wallet. It reads the aggregate mark-to-market of the trades the CSA governs \emph{in the real
ledger}, computes the required collateral (threshold, minimum transfer amount, eligible-collateral
and haircut rules), and emits margin-call moves to or from the collateral wallet. Margin
eligibility and haircut are contractual data in \textit{ProductTerms}; the exposure is the
observed valuation.

\noindent For a \emph{synthetic} managed account the netting set is the single real unit
$u_{\mathrm{TRS}}$, never the simulated constituents in the virtual ledger: reading those would
breach the isolation invariant (\S\ref{sec:trs}) and over-margin unmargined positions. The realm
tag and C4 realm-scoped reads enforce this. The margin exposure $N_k\cdot\mathrm{TR}_k$ is itself
an \emph{observation} of the virtual book lifted across the realm boundary as a number; that this
observation is unattested and unversioned is escalation~E2.

%==============================================================================
\section{TRS as Synthetic Managed Account}
\label{sec:trs}
%==============================================================================

A total return swap connects a virtual ledger $\mathcal{L}_v$ to real counterparty wallets in the
real ledger $\mathcal{L}_r$ through periodic cash settlement.

\begin{invariant}[Ledger isolation]
\label{inv:p7}
No move has one endpoint in $\mathcal{L}_v$ and the other in $\mathcal{L}_r$. (Invariant~P7.)
\end{invariant}

\noindent At reset $t_k$ the TRS contract in $\mathcal{L}_r$ \emph{observes} the virtual book's
performance and emits one real move of $\big|\mathrm{Payment}_k\big|$,
$\mathrm{Payment}_k=N_k\cdot\mathrm{TR}_k-N_k\cdot r_k\cdot\Delta t_k$, direction by sign:
\begin{lstlisting}
Move(from: w_payer, to: w_receiver, unit: USD,
     quantity: |Payment_k|, metadata: "TRS_NET_SETTLEMENT")
\end{lstlisting}
A minimal virtual ledger is not a second engine but a $\mathrm{realm}\in\{\text{real,virtual}\}$
tag enforcing Invariant~\ref{inv:p7}. The total return $\mathrm{TR}_k=(V^v_{t_k}-V^v_{t_{k-1}})/
V^v_{t_{k-1}}$ is partial: it requires $V^v_{t_{k-1}}>0$, an unstated precondition that must be a
typed failure, not a division by zero.

\begin{theorem}[Equivalence to the funded managed account]
The TRS settlement and the funded managed-account fee/PnL settlement are the \emph{same} ledger
operation: one net conservative cash move whose magnitude is a deterministic function of an
observed performance and a recorded baseline. The funded book plays the role of $\mathcal{L}_v$;
the periodic settlement is the TRS payment.
\end{theorem}
\noindent\textit{Proof sketch.} Both reduce to $\mathrm{Crystallise}(\text{magnitude},\text{from},
\text{to})$ with magnitude $=$ observed performance over the interval and direction by sign; both
reset a baseline to the post-settlement value; neither moves a position across a boundary.
The settlement \emph{primitive} is one; the \emph{products} differ, and that difference is a CDM
(reporting) distinction, not a ledger one: a CDM \texttt{TotalReturnSwap} is two-legged
($N\,\mathrm{TR}$ and $N\,r\,\Delta t$), so the report derives from the retained CDM
\texttt{BusinessEvent}, never from the netted move (\S\ref{sec:conformance}). \qed

\noindent\textbf{Price consistency.} $\mathcal{L}_v$ valuation and $\mathcal{L}_r$ settlement must
use one price vector $P_t$. Because Invariant~\ref{inv:p7} forbids a cross-ledger move, a price
divergence produces unexplained PnL with \emph{no} internal reconciliation path --- the one place
the ``no internal break'' guarantee can be silently voided. The framework asserts consistency in
prose but does not bind it; the structural fix (a shared, content-addressed snapshot) is
escalation~E2.

%==============================================================================
\section{Redemption and Wind-down}
\label{sec:redemption}
%==============================================================================

Wind-down strikes a final NAV, settles the final fee, returns capital, and retires the mandate.
Positions are liquidated against the market wallet; the cash is paid to the client's external
wallet; the mandate unit is returned to the manager. All closed-out rows are retained at zero
under the monotone carrier, preserving the audit trail, the final high-water mark, and the
quantities needed for tax reconstruction.
\begin{lstlisting}
Move(from: w_C, to: w_mkt, unit: <asset>, quantity: <held>)   -- liquidate
Move(from: w_mkt, to: w_C, unit: USD, quantity: <proceeds>)
Move(from: w_C, to: w_ext, unit: USD, quantity: <final NAV>)  -- capital-out
Move(from: w_C, to: w_M, unit: u_MA, quantity: 1)             -- retire mandate
\end{lstlisting}

%==============================================================================
\section{Balance-Sheet Substantiation at Account Level}
\label{sec:substantiation}
%==============================================================================

Each account's balance at any date is a deterministic projection of the move stream filtered to
its wallets; client statements, desk PnL, and regulatory position reports differ only in the
filter. Within the ledger there is no separate account-level \emph{quantity} record to reconcile
against. The move stream substantiates \emph{quantities and
provenance}; \emph{valuations} depend on the external price vector, and disclosures and legal
status require external evidence (consistent with the dual-valuation section of the main
specification). Net settlement is \emph{operational} netting: presenting a balance sheet net
requires, in addition, a legally enforceable right of set-off and the intent to settle net
(IAS~32.42 / ASC~210-20) --- a determination outside the ledger. The accounts'
assets~$=$~liabilities~$+$~equity reconciliation holds at the quantity level by P1 (the trial
balance balances by construction); the valued and classified presentation does not follow from P1
alone.

%==============================================================================
\section{Worked Example}
\label{sec:example}
%==============================================================================

One managed account $w_C$ under mandate $u_{\mathrm{MA}}$ issued by manager $w_M$; virtual market
wallet $w_{\text{mkt}}$; client external wallet $w_{\text{ext}}$. USD reference. Terms: management
fee $2\%/\text{yr}$, performance fee $20\%$ over a high-water mark (no hurdle), quarterly
($\Delta t=0.25$). Conventions: management fee on end-of-period NAV; performance fee on NAV net of
the management fee, above the mark; the mandate is non-valued. Every figure below was verified by
a decimal conservation check: $\sum_w w(u)=0$ holds for \emph{every} unit after \emph{every} move.

\begin{longtable}{p{4.6cm}p{7.2cm}r}
\toprule
\textbf{Step} & \textbf{Moves} ($\text{src}\to\text{dst}:\text{qty unit}$) & \textbf{NAV} \\
\midrule
A. Mandate issuance & $w_M\to w_C:1\ u_{\mathrm{MA}}$ & --- \\
B. Subscription & $w_{\text{ext}}\to w_C:1{,}000{,}000\ \text{USD}$ & $1{,}000{,}000$ \\
C. Trade (buy $5{,}000$ AAPL @100) & $w_C\to w_{\text{mkt}}:500{,}000\ \text{USD}$;\ \ $w_{\text{mkt}}\to w_C:5{,}000\ \text{AAPL}$ & $1{,}000{,}000$ \\
D. Q1 NAV (AAPL $\to 130$) & (price move; no ledger move) & $\mathbf{1{,}150{,}000}$ \\
E. Fee crystallisation & $w_C\to w_M:5{,}750\ \text{USD}$ (mgmt);\ \ $w_C\to w_M:28{,}850\ \text{USD}$ (perf) & $\mathbf{1{,}115{,}400}$ \\
F. Wind-down & $w_C\to w_{\text{mkt}}:5{,}000\ \text{AAPL}$;\ \ $w_{\text{mkt}}\to w_C:650{,}000\ \text{USD}$;\ \ $w_C\to w_{\text{ext}}:1{,}115{,}400\ \text{USD}$;\ \ $w_C\to w_M:1\ u_{\mathrm{MA}}$ & $0$ \\
\bottomrule
\end{longtable}

\noindent\textbf{Fee arithmetic (Step~E).} No capital flow occurred inside the period, so
$\mathrm{Perf}=(1{,}150{,}000-1{,}000{,}000)-0=150{,}000$. Management fee
$=1{,}150{,}000\times0.02\times0.25=5{,}750$. NAV net of management
$=1{,}144{,}250$. Performance fee $=0.20\times(1{,}144{,}250-1{,}000{,}000)=28{,}850$. NAV after
both fees $=1{,}115{,}400$. New high-water mark $=\max(1{,}000{,}000,\,1{,}115{,}400)=1{,}115{,}400$;
reset baseline $B_1=1{,}115{,}400$. Fee zero-sum:
$\sum_w\Delta_{\text{fee}}(\text{USD})=-34{,}600\,(w_C)+34{,}600\,(w_M)=0$.

\noindent\textbf{Closing identity (all wallets).} Final non-zero balances:
$w_M=+34{,}600$ (fee revenue), $w_{\text{mkt}}=-150{,}000$ (the gross gain paid out by the
market), $w_{\text{ext}}=+115{,}400$ (client net). Their sum is
$34{,}600-150{,}000+115{,}400=0$. The client received principal $1{,}000{,}000$ plus net gain
$115{,}400$; the gross gain $150{,}000$ equals net $115{,}400$ plus fees $34{,}600$. Per-unit
conservation holds across the whole chain: $u_{\mathrm{MA}}$ is issued $(+1,-1)$ then returned
$(0,0)$; USD nets to $0$ at $w_C$ and globally after every move; AAPL is acquired and liquidated
with $w_{\text{mkt}}$ as the mirror. Every position row is retained at zero under the monotone
carrier.

%==============================================================================
\section{Conformance and Flags}
\label{sec:conformance}
%==============================================================================

No conformance is asserted; each item is either demonstrated against a cited source or flagged for
the responsible owner. The CDM strata named below (\texttt{LegalAgreement}, \texttt{Trade},
\texttt{TotalReturnSwap}, \texttt{TradeState}, \texttt{BusinessEvent}) are identified at the type
level only; pinning them to a named CDM release is owed to the responsible CDM owner and is
\emph{flagged}, not asserted here.

\begin{itemize}[leftmargin=*]
\item \textbf{CDM mandate representation.} $u_{\mathrm{MA}}$ maps to the CDM
\texttt{LegalAgreement} stratum (the ISDA Create-digitised investment-management agreement:
parties, fee schedule, benchmark identity, methodology, limits), not the \texttt{Trade}/product
stratum. A discretionary mandate is a legal agreement, not a tradable contract.
\item \textbf{CDM TRS representation.} A CDM \texttt{TotalReturnSwap} is two-legged; the report is
the retained \texttt{BusinessEvent} (performance leg and financing leg
$N\,r\,\Delta t$), never the netted settlement move. The ledger ``same mechanism'' identity is a
settlement-primitive identity, \emph{flagged} as not a CDM/reporting identity.
\item \textbf{SFTR/EMIR reporting surface.} $u_{\mathrm{MA}}$ issuance is \emph{not} reportable: a
discretionary mandate is the MiFID~II Annex~I~\S A(4) portfolio-management service, not a Section~C
instrument, not an SFT, not a derivative. This discharges the open risk F5 of Addendum~A1. The
reportable artefact is never $u_{\mathrm{MA}}$ but the underlying book trade or $u_{\mathrm{TRS}}$,
each carrying its own UTI/UPI/LEI. Reportability is a per-relationship, regime-cited attribute on
the \emph{wrapped} unit's \textit{ProductTerms}; it is not inferable from the coarse unit-type key
(escalation~GOV, \S\ref{sec:escalations}).
\item \textbf{IFRS classification.} Economic NAV and economic PnL are derived here; accounting
classification of fair-value changes (FVTPL/FVOCI/amortised cost) and net presentation
(IAS~32.42) are outside scope and require external evidence. Accrued fees, modelled as conserved
moves rather than stored scalars, carry their own contra-entry and are auditable as double-entry.
\item \textbf{CDM \texttt{TradeState} alignment.} That the CDM \texttt{TradeState}-per-\texttt{Trade}
view and \textit{PositionState}$[w,u]$ stay replay-consistent (both projections of one minted
event) is asserted, not yet verified; it is the outstanding F6 deliverable (a Rosetta mapping
rerun against the three-map schema).
\end{itemize}

%==============================================================================
\section{Escalated Weaknesses}
\label{sec:escalations}
%==============================================================================

Of $140$ challenge questions, $105$ were answered and cleared by every touched lens; $35$ escalated
to framework-level weaknesses that \S 6 cannot resolve without changing the framework. They are
stated plainly; none was resolved by reshaping a question, narrowing its scope, or assuming away
its premise.

\begin{description}[leftmargin=0pt,style=nextline]
\item[\textbf{E1 --- Store-vs-derive economic scalars (an Addendum~A1 internal inconsistency).}]
A1 homes the high-water mark, entry NAV, accrued fee, and reset baseline as mutable
\textit{PositionState} scalars. Yet each is a \emph{fold} of the event log: the high-water mark is
a monotone max over recorded post-crystallisation NAVs, the accrued fee a fold over the rate
schedule and the NAV path, the baseline a fold over settlements. A1's own \texttt{first\_touch\_date}
ruling demotes exactly such a value to ``not state'' because caching it ``would create a
fold-inconsistency under back-dated corrections.'' Storing these scalars therefore re-creates the
internal-reconciliation surface the framework claims to abolish: under a back-dated price
correction the stored value desyncs from the recomputed fold. \emph{Direction of a fix:} hold these
as derived projections, not stored fields; their re-derivation requires recorded prices (E2). A
breaking mandate amendment makes this concrete --- carrying the high-water mark to the successor
unit must be re-derivation from the (corrected) stream, never a copy, on pain of a silent mark
reset.

\item[\textbf{E2 --- The observation surface is unattested (the linchpin).}]
The price vector $P_t$, the benchmark level, the TRS financing fix $r_k$, and intraday path
observations (barrier touch, realised vol) all strike real, irreversible cash. Each enters the
ledger as a logged event --- a settlement, a benchmark observation, a rate fix --- and
\textit{UnitStatus} only caches the latest such value for shared reading; the value is overwritten
in place only as a read cache, and every change to it is caused by one of those logged events.
\emph{Reproducing the number is therefore assured.} The exact price that struck a fee lives in the
immutable settlement event in the log, not only in the overwritten cache cell, so replaying the
events rebuilds it exactly: a valued balance sheet is as reproducible as the log itself --- as
reproducible as the events that produced it, never ``only as reproducible as a mutable cell.''
Path-independence (P10) holds across replay for the same reason, and because $\mathcal{L}_v$
valuation and $\mathcal{L}_r$ settlement both read the price from the same logged event, a fresh
replay reconstructs one price for both. Because the prices are in the log, E1's economic scalars
(high-water mark, baseline, accrued fee) are folds of recorded events, exactly as E1 requires.
\emph{What remains open is distinct from reproducibility: attestation.} The settlement payload
carries the price as a bare scalar with no attestation envelope --- provider key, source,
observation timestamp, fallback-chain-as-traversed, signature, snapshot content-address --- so the
number replays deterministically yet is not \emph{verifiable} to be the genuine settlement level:
determinism is not attestation. This attestation gap is recorded separately as the Nazarov
finding and is \emph{not} resolved here; it is orthogonal to where \textit{UnitStatus}'s value comes
from --- the source-of-truth question settled above.
The benchmark and $r_k$ moreover get weaker discipline than $P_t$, with no staleness, fallback, or
cross-source check, though the benchmark hurdle is client money. \emph{Direction of a fix
(attestation, the Nazarov finding):} the logged price, benchmark, and rate events must carry an
attested, content-addressed envelope, with cross-ledger consistency secured by sharing the snapshot
identity of the event rather than by after-the-fact reconciliation. It carries two governance
sub-items: the oracle-key lifecycle (generation, rotation, revocation --- on which the attestation,
not the replay, depends) is unspecified; and every trust assumption at the cash boundary is
currently ownerless, which is forbidden for a value-bearing input.

\item[\textbf{E3 --- The closed system is blind to counterparty legal identity.}]
P1 models an external counterparty as a virtual wallet, correct for \emph{quantity} closure but
erasing the counterparty's legal identity. Dual-sided EMIR/CFTC/SFTR reporting requires two
validated ISO-17442 LEIs, a reporting-counterparty determination, FC/NFC$(\pm)$ or SD/MSP
classification, and a paired UTI from the CPMI-IOSCO / EU Art~7 generation waterfall on the real
$u_{\mathrm{TRS}}$ --- none of which is a quantity fact, so none lives in the move stream. The only
candidate home, \textit{WalletRegistry}, is declared non-economic and non-versioned; a mutable LEI
would let an M\&A transfer or lapse retroactively rewrite every historical report on replay, so the
binding must be a time-indexed append-only sequence. Moreover a trade-repository \emph{pairing}
break is an \emph{external} fact: a closed ledger can produce our side of a report but cannot pair
it. \emph{Direction of a fix:} either promote counterparty LEI and classification to first-class,
time-versioned, replayable state bound to each externally-facing real wallet, or formally scope LEI
as an external reference-data dependency that the minting step joins against and records in the
event payload for deterministic replay.

\item[\textbf{E4 --- C4 capability scoping is asserted, not typed.}]
Segregation (\S\ref{sec:segregation}) and the no-cross-$(w,u_{\mathrm{MA}})$-read guarantee depend
on C4, but A1's reference accessor $\texttt{position}(w,u)$ takes no capability argument, so a
handler firing on one mandate can read another's state. This is answerable by adding a capability
parameter to the accessor (a section-level refinement); it is recorded here because it is the typed
expression on which Theorem~\ref{thm:seg} rests, and the reference signature does not yet carry it.

\item[\textbf{E5 --- No in-ledger solvency liveness.}]
Crystallising unrealised mark-to-market can drive the cash wallet negative; since negative is a
legal obligation, conservation alone cannot distinguish a \emph{funded} obligation from an
\emph{insolvent} overdraft. The ledger faithfully represents the obligation; funding is
settlement-layer state, and a client-money wallet carries a non-negativity precondition. The
residual --- a liveness invariant tying crystallisation to settlement finality --- is a boundary
property the closed ledger represents but does not guarantee.
\end{description}

\noindent\textbf{Reading of the escalations.} They are not failures of the workflow above, which is
sound as quantity algebra: every step is expressible as atomic $q>0$ moves with $\sum_w w(u)=0$
verified, and the worked example closes to the penny. They are the points at which the
\emph{performance and reporting} layer outruns the framework's current treatment of recorded
observations and external identity. E2 is the linchpin: the observation surface already enters as
logged events, so replay rebuilds every struck price exactly and E1's economic scalars are folds of
those events; what remains for E2 is \emph{attestation} --- an attested, content-addressed envelope
that makes the logged numbers verifiable, not merely reproducible (recorded separately as the
Nazarov finding). These weaknesses are the agenda for the next framework version.

\end{document}
